\documentclass[11pt,a4paper]{article}
\usepackage[T1]{fontenc}
\usepackage[utf8]{inputenc}
\usepackage{lmodern}
\usepackage[margin=26mm,headheight=14pt]{geometry}
\usepackage{amsmath,amssymb,amsthm,mathtools}
\usepackage{booktabs,longtable,array}
\usepackage{microtype}
\usepackage{xcolor}
\usepackage{enumitem}
\usepackage{fancyhdr}
\usepackage{xurl}
\usepackage[colorlinks=true,linkcolor=blue!45!black,citecolor=blue!45!black,urlcolor=blue!45!black]{hyperref}
\hypersetup{pdftitle={Resonance-Block Summation of Transseries},pdfauthor={Research report prepared for Vladimir Reshetnikov}}
\setlist{itemsep=3pt,topsep=5pt}
\setlength{\emergencystretch}{2em}
\numberwithin{equation}{section}
\newtheorem{theorem}{Theorem}[section]
\newtheorem{proposition}[theorem]{Proposition}
\newtheorem{lemma}[theorem]{Lemma}
\newtheorem{corollary}[theorem]{Corollary}
\theoremstyle{definition}
\newtheorem{definition}[theorem]{Definition}
\newtheorem{example}[theorem]{Example}
\theoremstyle{remark}
\newtheorem{remark}[theorem]{Remark}
\newcommand{\ii}{\mathrm i}
\newcommand{\e}{\mathrm e}
\newcommand{\dd}{\,\mathrm d}
\newcommand{\Res}{\operatorname{Res}}
\newcommand{\re}{\operatorname{Re}}
\newcommand{\im}{\operatorname{Im}}
\newcommand{\dist}{\operatorname{dist}}
\newcommand{\N}{\mathbb N}
\newcommand{\Z}{\mathbb Z}
\newcommand{\C}{\mathbb C}
\newcommand{\R}{\mathbb R}
\newcommand{\bw}{\boldsymbol\omega}
\newcommand{\sinc}{\operatorname{sinc}}
\newcommand{\exprel}{\operatorname{exprel}}
\newcommand{\norm}[1]{\left\lVert #1\right\rVert}
\pagestyle{fancy}
\fancyhf{}
\fancyhead[L]{\small Resonance-block summation of transseries}
\fancyhead[R]{\small Research report}
\fancyfoot[C]{\thepage}
\begin{document}
% ed. (2026-09-29): no page anchor on the title page, which otherwise shares its page number with the
% next page (pdfTeX "duplicate destination" warning); the same device the Stokes-transport package uses.
\hypersetup{pageanchor=false}
\begin{titlepage}
\vspace*{17mm}
\noindent{\large\scshape Transseries / Resurgence / Certified Inversion}\par
\vspace{12mm}
\noindent{\Huge\bfseries Resonance-Block\par Summation of Transseries}\par
\vspace{7mm}
\noindent{\Large Small divisors, coalescing Stokes actions,\par and nonlinear inversion}\par
\vspace{12mm}
\noindent Research report prepared for Vladimir Reshetnikov\par
\noindent 29 September 2026\par
\vspace{12mm}
\begin{abstract}
A countable Stokes expansion can have locally finite, finitely generated
positive action support and still fail to converge at every positive value
of the large variable. We establish this obstruction, and a constructive
repair, for the meromorphic Borel kernels
\[
 \widehat F_{\bw}(t)=\frac{t^d}{\prod_{j=1}^{d}\sin(\pi\omega_jt)},
 \qquad \omega_j>0.
\]
Nearby poles are grouped into blocks of at most $d$ nodes. Each block is a
confluent divided difference and admits a bound independent of every
internal pole separation. The block series equals the lateral Stokes jump,
converges normally throughout the right half-plane, and has an explicit
uniform action-cutoff error. For two incommensurate lattices, the ungrouped
series has the exact abscissa determined by exponential Diophantine
approximation; this abscissa can be infinite even though the original
series has exact Gevrey order one. We also prove a uniform coalescence
profile, characterize rational-exponential compression, and transport the
block sums through analytic inversion with an explicit geometric tail.
These results supply a concrete extension beyond a weighted individual-sector
hypothesis in the ProveIt repository. Complete conventional proofs and
reproducible checks are included; the claims are not Lean-formalized and
independent research priority is not asserted.
\end{abstract}
\vfill
\noindent\small Repository snapshot:
\texttt{0c973d8f5e7ef4550f4df1bf486d890037fd9f14}.\par
\noindent\small Source, PDF, verification programs, executed results, and build instructions
are distributed together. No repository files were modified.
\end{titlepage}
\hypersetup{pageanchor=true}
\tableofcontents
\newpage

\section{The question and the contribution}
\label{sec:intro}

% ed. (2026-09-29): scope note (title broader than the sine-product results); the abstract is unchanged.
\begin{quote}\small\textbf{Editorial note (ProveIt, 2026-09-29).} Scope. The title is broader than the results. Every theorem below concerns the sine-product Borel kernels \eqref{eq:kernel}, with a fixed finite number $d$ of period lattices; the one statement beyond them, the bounded-cluster criterion of Proposition~\ref{prop:abstract}, holds only under its stated hypotheses. No general resonance-block theory for transseries is proved. The ``resonances'' here are near-coincident Borel poles, not the equal-action resonances of the Stokes-transport and moving-fold packages of this series.\end{quote}

\subsection{A gap between formal support and analytic realization}
A familiar transseries calculation assigns an exponentially small term to
each Borel singularity. When the singularities lie at positive actions
$a$, the expected Stokes correction has the schematic form
\begin{equation}
 J(x)=2\pi\ii\sum_a p_a(x)\e^{-ax},\qquad x\longrightarrow+\infty.
 \label{eq:schematic}
\end{equation}
The polynomial $p_a$ records the principal part of a pole. Three different
questions must be kept separate. Is the support legitimate for formal
asymptotic bookkeeping? Does a function with these asymptotics exist? Does
the displayed infinite sum converge numerically at fixed $x$? In the model
studied here the first two answers are yes, while the third can be no for
every positive $x$.

The obstruction does not involve actions accumulating at a finite point.
Our actions form a finite union of positive lattices. Every bounded action
interval contains finitely many poles, and all actions belong to the
additive monoid generated by $1/\omega_1,\ldots,1/\omega_d$. The obstruction
is instead that distinct lattices can have extremely close nodes at
arbitrarily large actions. Their individual residues contain small
denominators. The corresponding finite group of residues cancels these
denominators, but an infinite expansion that estimates the residues
separately does not see the cancellation.

This observation leads to a concrete extension problem:
\begin{quote}
Can one construct a normally convergent Stokes representation, with
explicit uniform error bounds and a usable nonlinear inversion calculus,
without imposing a Diophantine condition on the periods?
\end{quote}
The answer is affirmative for every fixed finite number of lattices in
\eqref{eq:kernel} below. We also determine exactly when the ungrouped
representation converges in the two-lattice case.

\subsection{Relationship to the repository}
The repository's article \emph{Nonlinear Stokes Transport under Logarithmic
Inversion} proves a contour-based inverse formula and a countable extension
under a weighted summability assumption
\cite{repo-stokes}. Its label \texttt{cor:countable} requires the sum of
individual perturbation norms to be finite and sufficiently small. Its
finite-action collision proposition, \texttt{prop:collision}, treats
bounded amplitudes that combine by addition when actions coincide. These
are appropriate hypotheses and conclusions; this report does not correct
an error in them.

The present problem starts one step earlier. We obtain perturbations from
Borel residues whose amplitudes can individually diverge as two poles
approach one another. The weighted individual-residue norm can be infinite.
A finite resonance block, unlike either of its constituent residues, has
a controlled analytic norm. We prove the estimate that makes the
repository's analytic inversion method applicable after this change of
representation. The inverse coefficient formula itself remains classical.

% ed. (2026-09-29): cross-reference note (the Stokes-transport article's open question, answered in part here).
\begin{quote}\small\textbf{Editorial note (ProveIt, 2026-09-29).} This report answers, without citing it, the forward half of the Stokes-transport article's research question ``Countably many resurgent input poles'' (\path{Nonlinear_Stokes_Transport_Logarithmic_Inversion/article.tex} in \path{Analysis/Transseries/docs/series-and-transseries/}, the question following \texttt{cor:countable} in its list of research questions; this report inspected only lines 180--360 of that source). That question proposes normal convergence of the forward meromorphic Borel kernel as a starting hypothesis. For the kernels \eqref{eq:kernel}, Corollary~\ref{cor:counterexample} shows that normal convergence of the individual residues can fail at every positive argument, and Theorem~\ref{thm:block} restores it for resonance blocks with uniform constants. The resurgent half, an exact-core resurgent correction for the countable inverse, is not answered here (project~6 of Section~\ref{sec:questions}). Theorem~\ref{thm:inverse} is that article's \texttt{thm:transport} in core coordinates, as stated in Section~\ref{sec:inverse}.\end{quote}

The series-and-transseries inventory also contains recent studies of
finite-action folds, arithmetic action accumulation, and inverse harmonic
Stokes data \cite{repo-index}. The mechanism here is distinct: the number
of lattices is finite, their positive poles have no finite accumulation,
and the near-collisions occur at unbounded action. The source comparison is
limited to the inspected inventory and the cited Stokes article. It is not
a claim-by-claim audit of the repository's entire corpus.

\subsection{Classical ingredients and the boundary of the novelty claim}
Borel--Laplace summation and nonlinear closure are established subjects;
Sauzin's exposition is a suitable reference \cite{sauzin}. Divided
differences, including confluent nodes and their integral representations,
are classical \cite{deboor}. Small-divisor phenomena for quantum logarithms
and parameter-dependent functional equations are treated by Marmi and
Sauzin \cite{marmi-sauzin}. Meromorphic Borel transforms, residue expansions,
and Stokes jumps in quantum-dilogarithm problems are treated explicitly by
Garoufalidis and Kashaev \cite{garoufalidis-kashaev}.

Accordingly, neither cancellation of coalescing poles nor the general
existence of useful integral representations when some separated
expansions fail is claimed as a discovery. The contribution developed here
is a specific quantitative package: finite block cardinality, explicit
separation-independent bounds in arbitrary finite dimension, the exact
raw-series threshold for two lattices, a uniform collision law,
rational-exponential rigidity, and a blockwise inverse certificate. The
proofs establish the stated results. Whether every refinement is new to
the literature requires a broader priority review.

\subsection{Results at a glance}
The central theorem is Theorem~\ref{thm:block}: for a block $C$ with smallest
action $a_C$ and $m_C\le d$ pole occurrences, its contribution satisfies
\[
 |S_C(z)|\le K(1+a_C)^d\e^{-a_C\re z}
 \sum_{j=0}^{m_C-1}\frac{|z|^j}{j!\rho^{m_C-1-j}}.
\]
All constants are explicit and uniform on a compact positive period box.
No factor involves the separation of two nodes in the same block.
Theorem~\ref{thm:arithmetic} identifies the ungrouped abscissa as
\[
 \beta(\alpha)=\limsup_{n\to\infty}\frac1n
 \log\frac1{\dist(n\alpha,\Z)}.
\]
Theorems~\ref{thm:coalescence}, \ref{thm:rational}, and \ref{thm:inverse}
then give, respectively, the collision profile, the exact compression
criterion, and nonlinear transport. Section~\ref{sec:questions} records
further questions rather than treating these finite-lattice results as a
complete theory of countable-action resurgence.

\section{The Borel model and its lateral sums}
\label{sec:borel}

\subsection{Conventions}
Fix $d\ge1$ and positive real periods $\bw=(\omega_1,\ldots,\omega_d)$.
The starting point is
\begin{equation}
 \widehat F_{\bw}(t)=
 \frac{t^d}{\prod_{j=1}^{d}\sin(\pi\omega_jt)}.
 \label{eq:kernel}
\end{equation}
The value at zero is filled in by continuity:
\[
 \widehat F_{\bw}(0)=\frac1{\pi^d\omega_1\cdots\omega_d}.
\]
The function is even. Write its Taylor series as
\begin{equation}
 \widehat F_{\bw}(t)=\sum_{n\ge0}b_n t^n,
 \qquad
 \widetilde F_{\bw}(z)=\sum_{n\ge0}n!b_n z^{-n-1}.
 \label{eq:borel-convention}
\end{equation}
Thus our formal Borel map sends $z^{-n-1}$ to $t^n/n!$.
All odd $b_n$ vanish. The positive poles, counted with multiplicity, are
\begin{equation}
 \mathcal P_{\bw}=\{n/\omega_j:j=1,\ldots,d,\ n=1,2,\ldots\}.
 \label{eq:poles}
\end{equation}
Different occurrences at the same location represent a higher-order pole,
not separate simple residues. The numerator does not cancel any positive
or negative pole.

For $0<\theta<\pi/2$, define
\begin{equation}
 F_\pm(z)=\int_0^{\e^{\pm\ii\theta}\infty}
       \e^{-zt}\widehat F_{\bw}(t)\dd t.
 \label{eq:lateral}
\end{equation}
Initially the corresponding sufficient condition is
$\re(z\e^{\pm\ii\theta})>0$. Integrals with different admissible angles
on the same side agree on their common domains by contour deformation.
For any $z$ in the right half-plane, a sufficiently small positive
$\theta$ satisfies both inequalities. These choices patch to two
holomorphic functions on $\re z>0$. Throughout the article,
\begin{equation}
 J_{\bw}(z)=F_-(z)-F_+(z),
 \qquad S_{\bw}(z)=\frac{J_{\bw}(z)}{2\pi\ii}.
 \label{eq:jump-sign}
\end{equation}
This sign convention matters in every residue formula below.

\subsection{Uniform Gevrey estimates}
\begin{proposition}[Exact Gevrey order and off-ray realization]
\label{prop:gevrey}
The series $\widetilde F_{\bw}$ has Gevrey order one and no smaller
nonnegative Gevrey order. Its Borel transform is meromorphic with a locally
finite pole set. Fix
\begin{equation}
 0<m_0\le\omega_j\le M_0<\infty
 \quad(j=1,\ldots,d).
 \label{eq:box}
\end{equation}
There are constants $C,A>0$, uniform on this box, such that
$|n!b_n|\le CA^n n!$. For fixed $\theta\in(0,\pi/2)$ there are uniform
$C_\theta,A_\theta>0$ for which, for $x>0$ and $N\ge0$,
\begin{equation}
 \left|F_\pm(x)-\sum_{n=0}^{N-1}n!b_n x^{-n-1}\right|
 \le C_\theta A_\theta^N N!x^{-N-1}.
 \label{eq:gevrey-remainder}
\end{equation}
\end{proposition}
\begin{proof}
Choose $0<r<1/M_0$. On $|t|\le r$ the origin is removable and no other
zero of a denominator occurs. Compactness of the period box gives a
uniform bound there. Cauchy's inequality yields $|b_n|\le Cr^{-n}$.
For each fixed $\bw$, the nearest poles are at
$\pm1/\max_j\omega_j$ and they are genuine. The Borel Taylor radius is
therefore finite and equal to $1/\max_j\omega_j$.
If $|n!b_n|\le CB^n(n!)^s$ held with $s<1$, then
$|b_n|^{1/n}\le CB(n!)^{(s-1)/n}$ up to an inessential factor tending to
one. The right side tends to zero, making the Borel transform entire.
This is a contradiction.

On the ray $t=u\e^{\ii\theta}$ one has
\[
 |\sin(\pi\omega_jt)|^2
 =\sin^2(\pi\omega_ju\cos\theta)
  +\sinh^2(\pi\omega_ju\sin\theta).
\]
Consequently $\widehat F$ is uniformly bounded along that ray: it is
regular near zero and decreases exponentially, up to the factor $u^d$,
at infinity. The same holds below the axis.
Let $r_0=r/2$. Taylor's estimate on $u\le r_0$ and the coefficient bound
on $u\ge r_0$ imply a global ray estimate
\[
 \left|\widehat F(u\e^{\pm\ii\theta})
       -\sum_{n<N}b_n(u\e^{\pm\ii\theta})^n\right|
 \le C A^N u^N.
\]
For the second range, bound each polynomial term by
$Cu^N r_0^{-N}(r_0/r)^n$ and sum the resulting geometric series;
bound the function itself by its ray supremum. This proves the displayed
estimate with constants independent of $N$.
Integrating it against $\e^{-xu\cos\theta}$ gives
$CA^NN!/(x\cos\theta)^{N+1}$, which has the form
\eqref{eq:gevrey-remainder}. Integration of each Taylor monomial along
the ray gives $n!x^{-n-1}$. The same decay estimates justify the contour
patching in \eqref{eq:lateral}.
\end{proof}

\begin{remark}[What Gevrey regularity does not say]
The coefficient bound concerns the Borel germ near the origin. It does not
bound the residues at arbitrarily remote, nearly colliding singularities.
No inference from \eqref{eq:gevrey-remainder} to absolute convergence of
\eqref{eq:schematic} will be made. The counterexample in
Section~\ref{sec:arithmetic} shows why such an inference would be false.
\end{remark}

\section{A separation-independent resonance-block theorem}
\label{sec:blocks}

\subsection{The partition and explicit constants}
Continue to assume the compact box \eqref{eq:box}. Set
\begin{equation}
 \delta=\frac1{4dM_0},\qquad \rho=\frac\delta4.
 \label{eq:delta}
\end{equation}
List the positive pole occurrences in increasing order, keeping repeated
nodes. Join two consecutive occurrences precisely when their gap is less
than $\delta$. The connected components are called \emph{resonance
blocks}. This is a fixed deterministic partition once the periods and the
box are specified; it is not necessary to optimize the value of $\delta$.
For a block $C$, denote its ordered nodes by $a_1,\ldots,a_m$, its size by
$m=m_C$, and its minimum by $a_C$.

We use the deliberately conservative constants
\begin{align}
 B_0&=\max\left\{1,\frac2{\pi m_0},\frac1{m_0\delta}\right\},
 &L_0&=1+\frac1{4M_0},\notag\\
 K&=(L_0B_0)^d,
 &D_0&=d(\lceil M_0\rceil+2),\label{eq:constants}\\
 P_d(v)&=\sum_{m=1}^{d}\sum_{j=0}^{m-1}
          \frac{v^j}{j!\rho^{m-1-j}}.\notag
\end{align}
Sharper constants are possible. Their advantage here is that they require
no number-theoretic information about the periods.

\begin{lemma}[Finite block geometry]
\label{lem:geometry}
Every block contains at most $d$ occurrences and at most one occurrence
from each lattice. Its diameter is less than $(d-1)\delta$. Distinct blocks
are separated by at least $\delta$. The number of block minima in any
interval of length one is at most $D_0$.
\end{lemma}
\begin{proof}
Suppose $d+1$ consecutive occurrences were connected. Their span would be
less than $d\delta=1/(4M_0)$. Two of the occurrences would come from the
same one of the $d$ lattices. Distinct nodes in that lattice are separated
by at least $1/M_0$, a contradiction. The same argument excludes two
occurrences from one lattice in any block. The diameter and separation
claims follow from the construction.
A lattice of spacing $1/\omega_j$ has at most
$\lceil M_0\rceil+2$ nodes in any unit interval. Every block minimum is a
node of one of the lattices, so summing this bound proves the final claim.
\end{proof}

\subsection{Removing the nearby poles before estimating}
For each block define
\begin{equation}
 H_C(t)=\widehat F_{\bw}(t)\prod_{\nu=1}^{m_C}(t-a_\nu).
 \label{eq:HC}
\end{equation}
All factors associated with repeated nodes are retained in the product.
The point of \eqref{eq:HC} is that $H_C$, rather than any individual
residue, has a uniform local bound.

\begin{lemma}[Uniform regular-factor bounds]
\label{lem:regular}
The function $H_C$ is holomorphic in a neighborhood of the closed tube
\[
 \{t\in\C:\dist(t,[a_1,a_m])\le2\rho\}.
\]
Throughout this tube,
\begin{equation}
 |H_C(t)|\le K(1+a_C)^d.
 \label{eq:Hbound}
\end{equation}
For real $u\in[a_1,a_m]$ and every integer $k\ge0$,
\begin{equation}
 |H_C^{(k)}(u)|\le k!K(1+a_C)^d\rho^{-k}.
 \label{eq:derivativeH}
\end{equation}
\end{lemma}
\begin{proof}
By Lemma~\ref{lem:geometry}, each lattice either contributes one canceled
zero to the block or contributes none. For a contributing lattice, call
its node $a$. On the stated tube,
\[
 |t-a|\le(d-1)\delta+2\rho=(d-\tfrac12)\delta<\frac1{4M_0}.
\]
After removing the harmless sign $(-1)^n$,
$\sin(\pi\omega_jt)=\sin(\pi\omega_j(t-a))$. If $|w|\le\pi/4$, the
power series gives
\[
 \left|\frac{\sin w}{w}-1\right|
 \le\frac{\sinh|w|}{|w|}-1<\frac14.
\]
It follows, with a slightly relaxed constant, that
\[
 \left|\frac{t-a}{\sin(\pi\omega_jt)}\right|
 \le\frac2{\pi m_0}.
\]
The canceled quotient is holomorphic, including at $t=a$.

For a noncontributing lattice, no node lies in the block interval, and
all its positive nodes outside that interval are at least $\delta$ away.
Zero and negative nodes are farther away still, because the first positive
node of any lattice is at least $1/M_0$. Thus, throughout the tube,
\[
 \dist(\re t,\omega_j^{-1}\Z)\ge\delta-2\rho=\delta/2.
\]
For real $v$, $|\sin(\pi v)|\ge2\dist(v,\Z)$. Together with
$|\sin(u+\ii v)|\ge|\sin u|$, this yields
$|\sin(\pi\omega_jt)|\ge m_0\delta$.
Finally $|t|\le a_C+1/(4M_0)\le L_0(1+a_C)$.
Multiplying the bounds for the $d$ denominator factors and the numerator
proves \eqref{eq:Hbound}. Cauchy's estimate on a circle of radius $\rho$
centered at $u$ proves \eqref{eq:derivativeH}.
\end{proof}

\subsection{The cancellation identity}
For a function $f$ holomorphic near the nodes, its confluent divided
difference is characterized by
\begin{equation}
 [a_1,\ldots,a_m]f
 =\frac1{2\pi\ii}\int_\Gamma
      \frac{f(t)}{\prod_{\nu=1}^m(t-a_\nu)}\dd t,
 \label{eq:DDcontour}
\end{equation}
where $\Gamma$ surrounds the nodes and no singularity of $f$. At distinct
nodes this is
\begin{equation}
 [a_1,\ldots,a_m]f
 =\sum_{\nu=1}^m\frac{f(a_\nu)}{\prod_{\mu\ne\nu}(a_\nu-a_\mu)}.
 \label{eq:DDdistinct}
\end{equation}
At repeated nodes it is defined by continuous confluence; for $m$ equal
nodes it is $f^{(m-1)}(a)/(m-1)!$.

For real nodes we need the Hermite--Genocchi formula
\begin{equation}
 [a_1,\ldots,a_m]f
 =\int_{\Delta_{m-1}}f^{(m-1)}\!\left(\sum_{\nu=1}^m
                 \lambda_\nu a_\nu\right)\dd\lambda.
 \label{eq:HG}
\end{equation}
Here $\lambda_\nu\ge0$, $\sum\lambda_\nu=1$, and the coordinate measure
on the simplex has volume $1/(m-1)!$; it is not the induced Euclidean
surface measure. For completeness, the formula for $m=2$ is the
fundamental theorem of calculus after a linear change of variable.
For general $m$, integrate first along a simplex coordinate joining the
first and last nodes. The fundamental theorem gives the difference of the
two $(m-2)$-simplex integrals divided by $a_m-a_1$. Induction gives the
usual divided-difference recursion. Distinct nodes suffice for this
argument; continuity of the compact simplex integral gives repeated nodes.
Formula \eqref{eq:DDcontour} satisfies the same recursion, by subtracting
the two rational denominators, and the same initial case. This proves the
identities, including the confluent case. See also \cite{deboor}.

\begin{definition}[Block contribution]
The contribution of a block is the finite sum over its distinct locations
\begin{equation}
 S_C(z)=\sum_{a\in C\ \mathrm{distinct}}
            \Res_{t=a}\bigl(\e^{-zt}\widehat F_{\bw}(t)\bigr).
 \label{eq:blockdef}
\end{equation}
Equivalently, without any separation denominator,
\begin{equation}
 S_C(z)=[a_1,\ldots,a_m]\bigl(\e^{-zt}H_C(t)\bigr).
 \label{eq:blockDD}
\end{equation}
\end{definition}
Equation \eqref{eq:blockDD} follows immediately by substituting
\eqref{eq:HC} into \eqref{eq:DDcontour}. It is an exact identity, not an
asymptotic cancellation.

\subsection{Normal convergence, Stokes equality, and the tail}
\begin{theorem}[Uniform resonance-block summation]
\label{thm:block}
Under \eqref{eq:box}, for $z$ with $x=\re z>0$ and every block $C$,
\begin{equation}
 |S_C(z)|\le K(1+a_C)^d\e^{-a_Cx}
     \sum_{j=0}^{m_C-1}\frac{|z|^j}{j!\rho^{m_C-1-j}}.
 \label{eq:block-bound}
\end{equation}
The series $\sum_C S_C(z)$ converges absolutely and locally uniformly on
$x>0$, and
\begin{equation}
 F_-(z)-F_+(z)=2\pi\ii\sum_C S_C(z).
 \label{eq:Stokes-block}
\end{equation}
For $T\ge0$, retaining every block with $a_C<T$ gives the explicit bound
\begin{equation}
 \left|S_{\bw}(z)-\sum_{a_C<T}S_C(z)\right|
 \le
 \frac{KD_0d!(T+2)^dP_d(|z|)}{(1-\e^{-x})^{d+1}}\e^{-Tx}.
 \label{eq:tail}
\end{equation}
Every assertion and constant is uniform over the period box. The cutoff
is block-complete: a retained block may contain a node slightly above $T$.
\end{theorem}
\begin{proof}
Put $m=m_C$ and $f(t)=\e^{-zt}H_C(t)$. On the real block interval,
Leibniz's rule and \eqref{eq:derivativeH} give
\begin{align*}
 |f^{(m-1)}(u)|
 &\le K(1+a_C)^d\e^{-a_Cx}
 \sum_{j=0}^{m-1}\binom{m-1}{j}|z|^j(m-1-j)!\rho^{-(m-1-j)}\\
 &=(m-1)!K(1+a_C)^d\e^{-a_Cx}
   \sum_{j=0}^{m-1}\frac{|z|^j}{j!\rho^{m-1-j}}.
\end{align*}
The simplex volume in \eqref{eq:HG} proves \eqref{eq:block-bound}.
In particular the bound depends on the size of a block, not the distance
between any two of its nodes.

There are at most $D_0$ block minima in each interval
$[T+k,T+k+1)$. Hence
\begin{align*}
 \sum_{a_C\ge T}|S_C(z)|
 &\le KD_0P_d(|z|)\e^{-Tx}
       \sum_{k\ge0}(T+k+2)^d\e^{-kx}\\
 &\le KD_0P_d(|z|)(T+2)^d\e^{-Tx}
       \sum_{k\ge0}(k+1)^d\e^{-kx}.
\end{align*}
The elementary generating function for $(k+1)^d$ is an Eulerian
polynomial divided by $(1-q)^{d+1}$; the polynomial has nonnegative
coefficients summing to $d!$. Thus
$\sum_{k\ge0}(k+1)^d q^k\le d!/(1-q)^{d+1}$ for $0<q<1$.
Alternatively this inequality follows by expanding the powers in rising
factorials. This proves \eqref{eq:tail} and local uniform convergence.

It remains to identify the sum with the actual jump. Choose radii $R_
\ell\to\infty$ at midpoints of the gaps between successive blocks.
Each such radius is at least $\delta/2$ from every positive pole.
For fixed $z$ in the right half-plane choose a small $\theta>0$ with
$\re(z\e^{\ii\varphi})\ge c>0$ on $|\varphi|\le\theta$.
On the outer arc $t=R_\ell\e^{\ii\varphi}$, the denominator is uniformly
bounded away from zero when $|\im t|$ is small, and grows exponentially
when $|\im t|$ is not small. More explicitly, for a fixed small $h>0$,
if $|\im t|\le h$, then
$|R_\ell-\re t|=O(h^2/R_\ell)$; for large $\ell$ the real part stays at
least $\delta/4$ from every pole. If $|\im t|\ge h$, use
$|\sin(\pi\omega_jt)|\ge\sinh(\pi m_0h)$.
Thus $|\widehat F(t)|\le C R_\ell^d$ everywhere on these arcs.
Their integrals are $O(R_\ell^{d+1}\e^{-cR_\ell})$ and tend to zero.

Apply the residue theorem to the contour that travels out along the lower
ray, traverses the outer arc counterclockwise, and returns along the upper
ray. All enclosed poles form complete blocks. There is no pole at zero.
The limiting identity is \eqref{eq:Stokes-block}, with the sign prescribed
in \eqref{eq:jump-sign}. The same argument works locally uniformly in $z$.
\end{proof}

\begin{corollary}[Regularity across rational resonances]
\label{cor:parameter}
The total jump is real-analytic in the positive periods, including at
coincidences of any multiplicity at most $d$. Its values are independent
of which sufficiently small admissible block threshold is used.
\end{corollary}
\begin{proof}
For periods in a sufficiently small complex neighborhood of any fixed
positive vector, fixed upper and lower rays avoid every zero of the
corresponding denominators. Their imaginary parts remain bounded away
from zero relative to the ray parameter outside a neighborhood of the
origin. The integrals \eqref{eq:lateral} therefore depend holomorphically
on the periods, by locally uniform exponential domination. Their
difference has the same property. Theorem~\ref{thm:block} identifies every
admissible partition with this one difference. Individual blocks can
split or join when the partition changes; the total does not change.
\end{proof}

\begin{remark}[Not an unconditional rearrangement theorem]
When the raw residue series is divergent, \eqref{eq:Stokes-block} is not a
rearrangement of an already convergent series. It is a grouping prescription
justified by contour exhaustion and a new absolute block estimate.
Deleting one member of each close pair can destroy the conclusion. The
finite grouping is part of the analytic representation.
\end{remark}

\section{The exact arithmetic obstruction for two lattices}
\label{sec:arithmetic}

\subsection{Individual residues and the approximation exponent}
Specialize to $d=2$ and $(\omega_1,\omega_2)=(1,\alpha)$, where
$\alpha>0$ is irrational. All positive poles are simple. At the integer
poles and the second-lattice poles the residues of the Borel kernel are
\begin{equation}
 r_n=\frac{(-1)^n n^2}{\pi\sin(\pi\alpha n)},\qquad
 r'_m=\frac{(-1)^m(m/\alpha)^2}{\pi\alpha\sin(\pi m/\alpha)}.
 \label{eq:rawres}
\end{equation}
The ungrouped candidate for the normalized jump is
\begin{equation}
 \sum_{n\ge1}r_n\e^{-nz}
 +\sum_{m\ge1}r'_m\e^{-(m/\alpha)z}.
 \label{eq:rawseries}
\end{equation}
One may either ask that both series converge or order all their terms by
increasing action. The absolute convergence criterion is the same, and
the failure of the term test established below applies to either ordinary
interpretation.

Let $\|u\|=\dist(u,\Z)$, and define the exponential approximation exponent
\begin{equation}
 \beta(\alpha)=\limsup_{n\to\infty}\frac1n\log\frac1{\|n\alpha\|}
 \in[0,\infty].
 \label{eq:beta}
\end{equation}
This exponent measures exponential, not merely polynomial,
approximation. In particular, calling a number Liouville is not enough
to determine whether its $\beta$ is zero, finite positive, or infinite.

\begin{lemma}[Reciprocal covariance]
\label{lem:reciprocal}
For every positive irrational $\alpha$,
\begin{equation}
 \alpha\beta(1/\alpha)=\beta(\alpha).
 \label{eq:beta-reciprocal}
\end{equation}
\end{lemma}
\begin{proof}
Suppose first that $0<b<\beta(\alpha)$, so there are infinitely many $n$
with $\|n\alpha\|<\e^{-bn}$. Let $m$ be the nearest integer to
$n\alpha$. Then $m/n\to\alpha$, and for these sufficiently large $n$,
$n$ is the nearest integer to $m/\alpha$. Hence
\[
 \|m/\alpha\|=\frac{\|n\alpha\|}{\alpha},\qquad
 \limsup\frac1m\log\frac1{\|m/\alpha\|}\ge\frac b\alpha.
\]
Letting $b$ tend to $\beta(\alpha)$ proves
$\beta(1/\alpha)\ge\beta(\alpha)/\alpha$ whenever the latter is positive,
including the case of infinity. Interchanging $\alpha$ and $1/\alpha$
gives the reverse inequality. If one exponent is zero and the other were
positive, the same positive-exponent argument in the reverse direction
would give a contradiction. This covers all cases.
\end{proof}

\begin{theorem}[Exact ungrouped abscissa]
\label{thm:arithmetic}
For irrational $\alpha>0$, the raw residue series
\eqref{eq:rawseries} converges absolutely when $\re z>\beta(\alpha)$.
When $\re z<\beta(\alpha)$, its individual terms do not tend to zero.
Consequently its abscissa of absolute convergence is exactly
$\beta(\alpha)$, possibly infinite. No assertion is made here about
convergence on the boundary $\re z=\beta(\alpha)$.
In every case, the block representation of the same Stokes jump converges
normally on the whole half-plane $\re z>0$.
\end{theorem}
\begin{proof}
For real $u$,
\[
 2\|u\|\le|\sin(\pi u)|\le\pi\|u\|.
\]
The factors $n^2$ and constant denominators in \eqref{eq:rawres} do not
contribute to logarithmic growth after division by $n$. Therefore
\[
 \limsup_{n\to\infty}\frac1n\log|r_n|=\beta(\alpha).
\]
The root test proves absolute convergence of the first series for
$\re z>\beta(\alpha)$. If $\re z<\beta(\alpha)$, choose a number
strictly between them. Infinitely many coefficients grow faster than its
exponential, and the corresponding terms fail to tend to zero.
For the second lattice, the same root test in the index $m$ gives the
threshold $\re z>\alpha\beta(1/\alpha)$, which is identical by
Lemma~\ref{lem:reciprocal}. The terms of the first lattice already prove
failure below the common threshold, regardless of how the two lattices
are interleaved in an ordinary series. The final statement follows from
Theorem~\ref{thm:block} after choosing any compact positive box containing
$1$ and $\alpha$.
\end{proof}

\subsection{A benign irrational and an extreme irrational}
For $\alpha=\sqrt2$, let $m$ be nearest to $n\sqrt2$. Since
$2n^2-m^2$ is a nonzero integer,
\[
 \|n\sqrt2\|=|n\sqrt2-m|
 =\frac{|2n^2-m^2|}{n\sqrt2+m}\ge\frac c n
\]
for a fixed $c>0$ and all $n\ge1$. Hence $\beta(\sqrt2)=0$. The raw
residue sum converges absolutely for every $\re z>0$, although block
summation still improves conditioning near close poles.

In the opposite direction, define integers and a real number by
\begin{equation}
 N_1=2,\qquad N_{k+1}=10^{2N_k},\qquad
 \alpha_*=1+\sum_{k\ge1}10^{-N_k}.
 \label{eq:superliouville}
\end{equation}
This decimal has infinitely many nonzero digits and unbounded gaps
between them, so it is not eventually periodic and is irrational. It
lies strictly between one and two. Put $q_k=10^{N_k}$. The first $k$
terms of $q_k\alpha_*$ are integers, while the tail satisfies
\[
 0<\|q_k\alpha_*\|\le
   2\,10^{N_k-N_{k+1}}.
\]
For sufficiently large $k$ the tail is less than $1/2$, justifying the
nearest-integer notation. Since $N_{k+1}=q_k^2$,
\begin{equation}
 \frac1{q_k}\log\frac1{\|q_k\alpha_*\|}
 \ge\frac{(q_k^2-N_k)\log10-\log2}{q_k}
 \longrightarrow\infty.
 \label{eq:beta-infinity}
\end{equation}
Thus $\beta(\alpha_*)=\infty$.

\begin{corollary}[A simultaneous counterexample]
\label{cor:counterexample}
There is a formal series with all of the following properties:
\begin{enumerate}[label=(\roman*)]
\item exact Gevrey order one and upper and lower Borel--Laplace sums;
\item a meromorphic Borel transform having only simple poles, with no finite
accumulation of poles;
\item positive Stokes actions contained in a finitely generated, locally
finite additive monoid;
\item a raw individual-residue Stokes series that diverges at every
$z$ with $\re z>0$;
\item an absolutely normally convergent finite-block Stokes representation
on that same half-plane, with an explicit tail estimate.
\end{enumerate}
\end{corollary}
\begin{proof}
Take \eqref{eq:kernel} with $d=2$ and periods $(1,\alpha_*)$.
Proposition~\ref{prop:gevrey} gives (i) and (ii). The support lies in
$\{n+m/\alpha_*:n,m\in\Z_{\ge0}\}$; its positive generators ensure local
finiteness. Equations \eqref{eq:beta-infinity} and
Theorem~\ref{thm:arithmetic} prove (iv), while
Theorem~\ref{thm:block} proves (v).
\end{proof}
The corollary isolates a limitation of a possible general principle,
not of resurgence itself. Meromorphic continuation and Laplace
summation still work; the pointwise sum of the individually resolved
Stokes sectors is the operation that fails.

\section{Finite-cut transseries remain valid}
\label{sec:finitecut}
For a distinct positive pole $a$, let $m_a$ be its multiplicity and write
its principal part as
\begin{equation}
 \widehat F_{\bw}(t)=\sum_{\ell=1}^{m_a}
            \frac{c_{a,-\ell}}{(t-a)^\ell}+O(1).
 \label{eq:principal}
\end{equation}
Then its exact residue contribution is
\begin{equation}
 \Res_{t=a}(\e^{-zt}\widehat F_{\bw}(t))
 =\e^{-az}p_a(z),\qquad
 p_a(z)=\sum_{\ell=1}^{m_a}
        c_{a,-\ell}\frac{(-z)^{\ell-1}}{(\ell-1)!}.
 \label{eq:pa}
\end{equation}
The polynomial has degree exactly $m_a-1$, because the pole has that
multiplicity and its highest principal coefficient is nonzero.

\begin{theorem}[Finite-cut asymptotics and uniqueness]
\label{thm:finitecut}
Fix the positive periods. List their distinct positive poles increasingly
as $a^{(1)}<a^{(2)}<\cdots$. For each fixed $L\ge0$ and real
$x\to+\infty$,
\begin{equation}
 S_{\bw}(x)-\sum_{\ell=1}^{L}
       \e^{-a^{(\ell)}x}p_{a^{(\ell)}}(x)
 =\e^{-a^{(L+1)}x}p_{a^{(L+1)}}(x)
  +O\!\left(x^{d-1}\e^{-b_Lx}\right)
 \label{eq:finitecut}
\end{equation}
for some $b_L>a^{(L+1)}$. The coefficients in this polynomial-exponential
asymptotic expansion are unique. These statements hold even when the
infinite individual-residue series diverges at every fixed positive $x$.
\end{theorem}
\begin{proof}
Only finitely many blocks meet the finite set of nodes through
$a^{(L+1)}$. Include those entire blocks and, if necessary, the block
containing the next node as well. A finite collection of individual
residues accounts for all splitting of blocks caused by the desired
cutoff. Within that finite collection, every node beyond
$a^{(L+1)}$ has a strictly larger action. After it, the remaining
complete blocks have a minimum action bounded strictly above
$a^{(L+1)}$. Apply \eqref{eq:tail} at that latter fixed cutoff. Since
$P_d(x)=O(x^{d-1})$ and $(1-\e^{-x})^{-d-1}$ is bounded for $x\ge1$,
this proves \eqref{eq:finitecut} for a suitable $b_L$.

For uniqueness, suppose two such expansions on locally finite positive
action supports disagree. Their union is still locally finite and has a
smallest action at which the polynomial coefficients differ. Subtract
the expansions and multiply by the exponential of that action times $x$.
The nonzero difference polynomial would be exponentially small times a
fixed power of $x$, which is impossible as $x\to+\infty$. If the two
supports include an action zero, the same argument starts there. Repeating
at each finite cutoff proves uniqueness.
\end{proof}

The constants in \eqref{eq:finitecut} may be enormous for a poorly
approximable cutoff choice, because a fixed near-collision may be split.
They are finite, which is all an asymptotic expansion at a fixed cutoff
requires. In contrast, \eqref{eq:tail} supplies useful uniform constants
when whole blocks are retained. This is the distinction between a formal
or Poincar\'e-style expansion and a numerical evaluation scheme.

One should also distinguish finite-generation of the action monoid from
finite-generation of a \emph{function field}. Formal products have locally
finite action support because sums of finitely many positive generators
have only finitely many exponent tuples below any bound. This support
fact alone does not give any analytic norm estimate for a product or an
inverse. The block norms in Section~\ref{sec:inverse} supply that missing
information.

\section{Coalescence and its universal scaled profile}
\label{sec:coalescence}

\subsection{An entire divided-difference profile}
For $m\ge1$, define
\begin{equation}
 \Phi_m(\lambda_1,\ldots,\lambda_m)
 =[\lambda_1,\ldots,\lambda_m]\e^{-u}.
 \label{eq:Phi}
\end{equation}
At distinct nodes,
\begin{equation}
 \Phi_m(\lambda_1,\ldots,\lambda_m)
 =\sum_{j=1}^m\frac{\e^{-\lambda_j}}
                 {\prod_{k\ne j}(\lambda_j-\lambda_k)}.
 \label{eq:Phi-simple}
\end{equation}
The apparent singularities in \eqref{eq:Phi-simple} are removable. The
representation
\begin{equation}
 \Phi_m(\lambda)=(-1)^{m-1}
       \int_{\Delta_{m-1}}\e^{-\sum\lambda_jv_j}\dd v
 \label{eq:Phi-integral}
\end{equation}
makes it entire and symmetric in its $m$ complex arguments. In particular,
\begin{equation}
 \Phi_m(\lambda,\ldots,\lambda)
   =\frac{(-1)^{m-1}}{(m-1)!}\e^{-\lambda}.
 \label{eq:Phi-confluent}
\end{equation}
For real arguments its sign is $(-1)^{m-1}$. These properties are useful
both for asymptotic analysis and for stable numerical evaluation.

\begin{theorem}[Uniform $x^{-1}$ action-coalescence law]
\label{thm:coalescence}
Let $x\to+\infty$, and suppose an $m$-node cluster has nodes
\begin{equation}
 a_j(x)=a+\frac{\lambda_j}{x}+O(x^{-2}),\qquad j=1,\ldots,m,
 \label{eq:movingnodes}
\end{equation}
uniformly for $\lambda$ in a fixed compact subset of $\R^m$, with repeated
nodes allowed. Suppose its regular factor $H_x$ is holomorphic on a
fixed neighborhood of $a$, its derivatives through order $m-1$ are
uniformly bounded there, and $H_x(a)=H_0+O(x^{-1})$, uniformly in those
parameters. Its residue block then satisfies
\begin{equation}
 S_C(x)=\e^{-ax}x^{m-1}
       \left(H_0\Phi_m(\lambda_1,\ldots,\lambda_m)+O(x^{-1})\right),
 \label{eq:coalescence}
\end{equation}
uniformly on the compact parameter set. The conclusion does not require
a positive lower bound on $|\lambda_j-\lambda_k|$.
\end{theorem}
\begin{proof}
Let $\mu_j=x(a_j(x)-a)=\lambda_j+O(x^{-1})$. The affine change of variable
$u=x(t-a)$ in a divided difference gives exactly
\[
 S_C(x)=\e^{-ax}x^{m-1}
 [\mu_1,\ldots,\mu_m]\bigl(\e^{-u}H_x(a+u/x)\bigr).
\]
On any fixed compact set containing the convex hulls of the $\mu_j$,
the function $H_x(a+u/x)-H_0$ is $O(x^{-1})$, as are its derivatives
in $u$ of orders up to $m-1$. This follows from the assumed derivative
bounds and the chain rule; for order zero also use
$H_x(a)-H_0=O(x^{-1})$.
The integral formula \eqref{eq:HG} therefore changes the last divided
difference by $O(x^{-1})$ when this function is replaced by $H_0$.
Formula \eqref{eq:Phi-integral} is uniformly Lipschitz on compact sets,
so replacing $\mu_j$ by $\lambda_j$ creates another $O(x^{-1})$ error.
No estimate has divided by a node separation. This proves the result.
\end{proof}

Theorem~\ref{thm:coalescence} explains the polynomial factor generated by a
multiple Borel pole. When $m$ actions approach one another on the
$x^{-1}$ scale, the combined contribution is of order
$x^{m-1}\e^{-ax}$, with a finite profile interpolating all collision
patterns. Exact coincidence is not a singular limit of the block.

For the sine-product model, the regular-factor hypotheses follow locally
by canceling the selected lattice zeros before varying the periods.
One must specify the collision under study and keep any other zeros out
of the chosen neighborhood; arbitrary additional collisions would change
$m$. Lemma~\ref{lem:regular} supplies uniform bounds for full blocks.

\subsection{Two coincident lattices}
For $d=2$, suppose $a=n=m/\alpha$ is a common positive node. Expanding
around $t=a+u$ gives
\[
 \widehat F_{(1,\alpha)}(a+u)
 =\frac{(-1)^{n+m}(a+u)^2}
        {\sin(\pi u)\sin(\pi\alpha u)}.
\]
The coefficient of $u^{-1}$ after multiplication by $\e^{-x(a+u)}$ is
\begin{equation}
 S_{\{a,a\}}(x)=
 \frac{(-1)^{n+m}}{\pi^2\alpha}
       \e^{-ax}(2a-a^2x).
 \label{eq:double-residue}
\end{equation}
At $\alpha=1$ every positive node is double. Summing the elementary
geometric derivatives yields the exact expression
\begin{equation}
 S_{(1,1)}(x)=\frac1{\pi^2}
 \left\{\frac{2q}{(1-q)^2}
       -\frac{xq(1+q)}{(1-q)^3}\right\},\qquad q=\e^{-x}.
 \label{eq:alpha1}
\end{equation}
It is valid for every $\re x>0$.

As a moving example take $\alpha=1+\lambda/x$ and the first two poles
$1$ and $1/\alpha$. Their scaled offsets are $0$ and $-\lambda$, and
$H_0=1/\pi^2$. Therefore
\begin{equation}
 \frac{\pi^2\e^x}{x}S_{\{1,1/\alpha\}}(x)
 \longrightarrow-\frac{\e^\lambda-1}{\lambda},
 \label{eq:two-profile}
\end{equation}
with value $-1$ at $\lambda=0$. This is a uniform compact-$\lambda$
limit, not just a statement for distinct limiting actions.
At $\lambda=0$, the finite-$x$ expression on the left is exactly
$-1+2/x$.

\subsection{A triple-pole check}
For $(\omega_1,\omega_2,\omega_3)=(1,1,1)$,
\[
 \frac1{\sin^3(\pi u)}
 =\frac1{\pi^3u^3}\left(1+\frac{\pi^2u^2}{2}+O(u^4)\right).
\]
Extracting the residue at the integer $n$ gives
\begin{equation}
 S_{\{n,n,n\}}(x)
 =\frac{(-1)^n\e^{-nx}}{\pi^3}
 \left(3n-3n^2x+\frac{n^3}{2}(x^2+\pi^2)\right).
 \label{eq:triple-residue}
\end{equation}
The leading term has the $x^2\e^{-nx}$ behavior predicted by
\eqref{eq:coalescence} and
$\Phi_3(0,0,0)=1/2$. This example also shows why limiting residues must
be taken as higher-order residues rather than by substituting a resonant
parameter into either simple-pole formula.

\section{Rational-exponential compression and its converse}
\label{sec:rational}

\begin{definition}
A function on a positive half-line has a \emph{finite single-base
rational-exponential representation} if for some $h>0$, some integer
$N\ge0$, and rational functions $R_k$ holomorphic at the origin, it can
be written
\begin{equation}
 S(x)=\sum_{k=0}^{N}x^k R_k(\e^{-hx}).
 \label{eq:rationalform}
\end{equation}
This is a specific compression class. It does not include arbitrary
rational functions of several independent exponential bases, nor every
special-function representation.
\end{definition}

\begin{theorem}[Exact two-lattice compression criterion]
\label{thm:rational}
For the two-lattice kernel $\widehat F_{(1,\alpha)}$, the normalized jump
$S_{(1,\alpha)}$ has a representation \eqref{eq:rationalform} if and only
if $\alpha$ is rational. If $\alpha=p/q$ in lowest terms with $p,q>0$,
one may take $h=1/p$ and $N\le1$.
More generally, for $d$ commensurate positive periods one may take
$N\le d-1$ with a suitable $h$.
\end{theorem}
\begin{proof}
Suppose first that the lattice spacings are commensurate. There are
$h>0$ and positive integers $r_j$ with $1/\omega_j=hr_j$. All nodes have
the form $h\ell$, and the multiplicity at that node is the number of
indices for which $r_j$ divides $\ell$. This pattern is periodic in
$\ell$.
Expand each denominator around $t=h\ell+u$. Whether its sine vanishes,
and every coefficient needed from its Laurent or Taylor expansion,
depend periodically on $\ell$, with a common period dividing a multiple
of $2\operatorname{lcm}(r_1,\ldots,r_d)$. The numerator $(h\ell+u)^d$
is polynomial in $\ell$. The residue polynomial \eqref{eq:pa} thus has
the form
\[
 p_{h\ell}(x)=\sum_{k=0}^{d-1}x^k A_k(\ell),
\]
where each $A_k$ is polynomial in $\ell$ separately on each of finitely
many residue classes. Terms at locations that are not poles are set to
zero. For $Q=\e^{-hx}$, every series
$\sum_{\ell\ge1}A_k(\ell)Q^\ell$ is rational in $Q$: split into residue
classes and apply powers of $Q\partial_Q$ to geometric series.
The coefficients have only polynomial growth, so these sums converge
absolutely when $|Q|<1$. This proves \eqref{eq:rationalform} with
$N\le d-1$ and rational functions holomorphic at zero.
For $(1,p/q)$ the spacings are $1$ and $q/p$, so $h=1/p$ works and
$d-1=1$.

Conversely suppose $\alpha$ is irrational and
\eqref{eq:rationalform} holds. Expanding the rational functions at zero
gives a convergent polynomial-exponential expansion supported on
$\{h\ell:\ell\ge0\}$. The zero-action polynomial vanishes because
$S_{(1,\alpha)}(x)$ is exponentially small. Every pole in the actual
kernel is simple and has a nonzero residue. By
Theorem~\ref{thm:finitecut}, its asymptotic action expansion is unique.
In particular the actions $1$ and $1/\alpha$ must both be positive integer
multiples of $h$. Their ratio would then be rational, contradicting the
irrationality of $\alpha$. This argument uses finite-cut asymptotics,
not convergence of the raw infinite series, so it also applies when
$\beta(\alpha)=\infty$.
\end{proof}

The theorem gives a precise contrast. The total jump varies analytically
through rational periods by Corollary~\ref{cor:parameter}, yet admitting a
finite single-base representation is an arithmetic property attained
exactly on rational two-lattice ratios. Smooth dependence of a function
on a parameter does not imply uniform finite symbolic compressibility.

\section{Nonlinear inversion in the block norm}
\label{sec:inverse}

\subsection{The analytic space that survives the small divisors}
The appropriate norm is taken after finite cancellation. Let $D$ be a
closed disk in the right half-plane. Theorem~\ref{thm:block} gives
\begin{equation}
 \sum_C\sup_{z\in D}|S_C(z)|<\infty.
 \label{eq:normalnorm}
\end{equation}
The corresponding sum over individual residues need not be finite.
For the example $\alpha_*$, it is infinite on every such disk with
nonempty interior, because even evaluation at any one point produces
terms that do not tend to zero.

We shall call \eqref{eq:normalnorm} the \emph{block norm}. This terminology
refers to a chosen representation in a Banach space of holomorphic
functions on a disk. It does not assert the construction of a new full
field of formal transseries. In particular, different decompositions can
have different norms even when their sums are the same function.

\subsection{A branch certificate and all-order inverse transport}
\begin{theorem}[Blockwise inverse transport and two truncation errors]
\label{thm:inverse}
Let $\phi$ be univalent and holomorphic on a neighborhood of
$\overline D(y_0,R)$, let $z_0=\phi(y_0)$, and let $G=\phi^{-1}$ be the
corresponding local core. Suppose
\begin{equation}
 j(y)=J(\phi(y))=\sum_C j_C(y),\qquad
 B=\sum_C\sup_{|y-y_0|\le R}|j_C(y)|<\infty,
 \label{eq:Bnorm}
\end{equation}
where the $j_C$ are holomorphic on a common neighborhood of the closed
disk, and the sum is holomorphic there or on an open neighborhood of the
disk after an arbitrarily small reduction of $R$. Put
\[
 M=\sup_{|y-y_0|\le R}|\phi'(y)|.
\]
If $r=|\lambda|B/R<1$, the equation
\begin{equation}
 G(z_\lambda)+\lambda J(z_\lambda)=y_0
 \label{eq:inverse-equation}
\end{equation}
has a unique solution in $\phi(D(y_0,R))$. It depends holomorphically on
$\lambda$ and satisfies
\begin{equation}
 z_\lambda-z_0
 =\sum_{n\ge1}\frac{(-\lambda)^n}{n!}
   \left.\partial_y^{n-1}\bigl(\phi'(y)j(y)^n\bigr)\right|_{y=y_0}.
 \label{eq:inverse-series}
\end{equation}
Expanding each $j^n$ over ordered $n$-tuples of blocks is absolutely
convergent. The sum of absolute values of all degree-$n$ contributions
is at most $MRr^n/n$, and the tail beyond degree $N\ge0$ is bounded by
\begin{equation}
 \frac{MR}{N+1}\frac{r^{N+1}}{1-r}.
 \label{eq:inverse-tail}
\end{equation}

For a subfamily of retained blocks let $j_T$ be its sum, and assume
\begin{equation}
 \Delta_T=\sum_{C\ \mathrm{omitted}}
        \sup_{|y-y_0|\le R}|j_C(y)|.
 \label{eq:Delta}
\end{equation}
The full inverse and the inverse formed with $j_T$ differ by at most
\begin{equation}
 \frac{M|\lambda|\Delta_T}{1-r}.
 \label{eq:action-inverse-tail}
\end{equation}
Thus truncating both the blocks and the amplitude degree gives a total
error bounded by the sum of \eqref{eq:action-inverse-tail} and
\eqref{eq:inverse-tail}. If $B=0$, all statements have the trivial
interpretation $z_\lambda=z_0$.
\end{theorem}
\begin{proof}
Use the core coordinate $y=G(z)$. The equation becomes
$y+\lambda j(y)=y_0$. On $|y-y_0|=R$, its perturbation has modulus at
most $|\lambda|B<R$. Rouch\'e's theorem gives exactly one zero in the
disk, counted with multiplicity. It is therefore simple. The implicit
function theorem and uniqueness give a single holomorphic branch in
$|\lambda|<R/B$.

The argument principle, applied with the holomorphic weight $\phi(y)$,
expresses the difference of the roots after applying $\phi$ as
\[
 \frac1{2\pi\ii}\int_{|y-y_0|=R}\phi(y)
 \left(\frac{1+\lambda j'(y)}{y-y_0+\lambda j(y)}
              -\frac1{y-y_0}\right)\dd y.
\]
The logarithm $\log(1+\lambda j(y)/(y-y_0))$ is single-valued on a thin
annulus containing the contour, by its uniformly convergent power series.
Integration by parts changes the preceding expression to
\begin{equation}
 -\frac1{2\pi\ii}\int_{|y-y_0|=R}
        \phi'(y)\log\left(1+\frac{\lambda j(y)}{y-y_0}\right)\dd y.
 \label{eq:inverse-contour}
\end{equation}
Expanding the logarithm and applying Cauchy's formula proves
\eqref{eq:inverse-series}.

For the absolute estimate, retain the contour form of each homogeneous
term. The sum of the moduli of all ordered products of block factors is
bounded on the contour by
$(\sum_C\|j_C\|)^n=B^n$. The contour length is $2\pi R$, the core
denominator has modulus $R^n$, and the logarithmic coefficient is $1/n$.
Consequently the total degree-$n$ bound is $MRr^n/n$. It justifies every
interchange of sums, derivatives, and contour integrals. Summing for
$n\ge N+1$ and using $1/n\le1/(N+1)$ proves
\eqref{eq:inverse-tail}.

For the second statement use \eqref{eq:inverse-contour} for $j$ and
$j_T$. Along the segment between their two contour values the arguments
of the logarithm remain in $|w-1|\le r$. Hence the difference of the two
logarithms has modulus at most
$|\lambda|\Delta_T/(R(1-r))$. Multiplying by the contour length and
$M/(2\pi)$ gives \eqref{eq:action-inverse-tail}. Both equations have a
unique root by the same strict boundary bound. This finishes the proof.
\end{proof}

The contour proof is a form of Lagrange inversion, and the corresponding
weighted-family result already appears in the repository
\cite{repo-stokes}. We included the proof to make the logical dependence
and constants explicit. The new input is the finite-block estimate that
verifies the norm hypothesis in examples where the individual-sector
version fails.

\subsection{Explicit certificates for the sine-product family}
Take $J=J_{\bw}$ from \eqref{eq:jump-sign} and
$j_C(y)=2\pi\ii S_C(\phi(y))$. Suppose the core image satisfies
\begin{equation}
 \re\phi(y)\ge x_*>0,\qquad |\phi(y)|\le V
 \quad (|y-y_0|\le R).
 \label{eq:core-image}
\end{equation}
Since every block minimum is at least $1/M_0$, the proof of
Theorem~\ref{thm:block} gives the directly evaluable certificate
\begin{equation}
 B\le
 \frac{2\pi KD_0d!(2+M_0^{-1})^dP_d(V)}
      {(1-\e^{-x_*})^{d+1}}\,
       \e^{-x_*/M_0}.
 \label{eq:Bcertificate}
\end{equation}
For the block-complete cutoff $a_C<T$,
\begin{equation}
 \Delta_T\le
 \frac{2\pi KD_0d!(T+2)^dP_d(V)}
      {(1-\e^{-x_*})^{d+1}}\e^{-Tx_*}.
 \label{eq:Deltacertificate}
\end{equation}
These are bounds on sums of suprema, not merely bounds on the absolute
value of the total perturbation. This distinction is essential to
absolute mixed-block convergence.

Equations \eqref{eq:Bcertificate}, \eqref{eq:inverse-tail}, and
\eqref{eq:Deltacertificate} provide two independent accuracy controls.
Increasing $T$ reduces the unresolved action tail; increasing $N$ reduces
the nonlinear amplitude tail. Neither control requires a lower bound on
an internal action separation or an estimate of $\beta(\alpha)$.

\begin{corollary}[A genuine extension beyond individual-sector inversion]
\label{cor:inverse-obstruction}
For the irrational $\alpha_*$ in \eqref{eq:superliouville}, the
block-indexed inverse expansion of Theorem~\ref{thm:inverse} is applicable
whenever its branch certificate holds. In contrast, an inverse expansion
resolved into the raw individual residues fails the term test already
in its first homogeneous degree at every base point $z_0$ with
$\re z_0>0$.
\end{corollary}
\begin{proof}
The block norm is finite by \eqref{eq:Bcertificate}. The degree-one
coefficient is
\begin{equation}
 -\lambda\phi'(y_0)J(z_0).
 \label{eq:first-inverse}
\end{equation}
Expanding $J(z_0)$ into individual residues yields terms equal to the
raw terms in Theorem~\ref{thm:arithmetic}, multiplied by the nonzero
constant $-2\pi\ii\lambda\phi'(y_0)$. For nonzero $\lambda$ they fail
to tend to zero. In block form the same degree-one coefficient is an
absolutely convergent sum. Higher-order rearrangements cannot repair
failure of absolute convergence at the first homogeneous degree.
\end{proof}

\subsection{Formal action preservation with analytic realization}
Let $\mathcal A$ be the distinct positive pole set, let
$a_*=\min\mathcal A$, and put $q_a(x)=2\pi\ii p_a(x)$, with $p_a$ from
\eqref{eq:pa}. The formally resolved jump is
$\widetilde J(x)=\sum_{a\in\mathcal A}q_a(x)\e^{-ax}$.

\begin{theorem}[An explicit finite formula at each inverse action]
\label{thm:formal-inverse}
For the identity core, the formal inverse of
$x\mapsto x+\lambda\widetilde J(x)$ is well-defined on the additive action
monoid $\langle\mathcal A\rangle$. Its correction has coefficients
\begin{align}
 \widetilde Z_\lambda(x)-x
   &=\sum_{A\in\langle\mathcal A\rangle\setminus\{0\}}
                   Q_A(x;\lambda)\e^{-Ax},\label{eq:formal-inverse}\\
 Q_A(x;\lambda)
   &=\sum_{1\le n\le A/a_*}\frac{(-\lambda)^n}{n!}
       \sum_{\substack{a_1,\ldots,a_n\in\mathcal A\\
                       a_1+\cdots+a_n=A}}
       (\partial_x-A)^{n-1}\prod_{j=1}^{n}q_{a_j}(x).
       \label{eq:action-coefficient}
\end{align}
Every sum defining $Q_A$ is finite. Thus $Q_A$ is polynomial in $x$ and
$\lambda$, and no new action outside the additive input monoid is needed.
For fixed $\lambda\in\C$, the actual inverse near large positive $x$ of
$z\mapsto z+\lambda J_{\bw}(z)$ has this expansion at every fixed finite
action cutoff. Its block-indexed amplitude series is convergent for all
sufficiently large $x$, even in the counterexample of
Corollary~\ref{cor:counterexample}.
\end{theorem}
\begin{proof}
The monoid generated by the $d$ positive spacings is locally finite, and
$\mathcal A$ is contained in it. A product of $n$ positive-action terms
has action at least $na_*$. Hence only $n\le A/a_*$ and finitely many
action tuples can contribute at action $A$. Differentiation does not
change an action and satisfies
\[
 \partial_x^{n-1}\bigl(\e^{-Ax}P(x)\bigr)
   =\e^{-Ax}(\partial_x-A)^{n-1}P(x).
\]
The identity-core case of the Lagrange formula therefore gives
\eqref{eq:action-coefficient}. It can also be verified in each finite
action quotient: substitution by a positive-action displacement raises
its filtration, so the usual recursive inverse is unique, and the finite
Lagrange identities hold at each such quotient. This proves formal
existence and uniqueness without an analytic convergence assumption on
$\widetilde J$.

For the analytic assertion take a fixed-radius disk about real $x$ and
apply Theorem~\ref{thm:inverse} with $\phi(y)=y$. Its block norm is
$O(x^{d-1}\e^{-a_*x})$; a constant factor from the disk radius is harmless.
It tends to zero, so every fixed $\lambda$ eventually satisfies the branch
certificate. Fix an action bound $A_0$. Choose $N$ with
$(N+1)a_*>A_0$. The amplitude tail \eqref{eq:inverse-tail} is then bounded
by a fixed power of $x$ times $\e^{-(N+1)a_*x}$ and is smaller than every
retained term at action at most $A_0$.
For each of the finitely many degrees $n\le N$, retain finitely many
whole jump blocks through an action strictly above $A_0$. The omitted
jump tail is exponentially smaller, uniformly on a slightly larger
fixed-radius disk, by \eqref{eq:tail}. Cauchy's formula preserves this
exponential margin after the finitely many derivatives. The retained
finite products give exactly \eqref{eq:action-coefficient} through
$A_0$. This proves finite-cut realization. It does not assert convergence
of the infinite raw action sum at fixed $x$.
\end{proof}

\subsection{Actual lateral maps}
For a fixed scalar $\varepsilon$, consider
\begin{equation}
 T_\pm(z)=z+\varepsilon F_\pm(z).
 \label{eq:actual-maps}
\end{equation}
On a fixed-width neighborhood of a sufficiently large positive real
point, the ray estimates give $F_\pm(z)=O(z^{-1})$ and
$F_\pm'(z)=O(z^{-2})$. These follow either by differentiating the
convergent ray integral or by its first Taylor-remainder estimate on a
slightly larger disk. Thus $T_+$ is a small derivative perturbation of
the identity and has a local inverse core $\phi$ on a disk of fixed
positive radius, after moving the base point sufficiently far right.

Set $G=T_+$ and $\lambda=\varepsilon$ in
Theorem~\ref{thm:inverse}. Since $T_-=T_++\varepsilon J$, its solution is
$T_-^{-1}(y_0)$ on the certified branch. This realizes the inverse Stokes
transport as an actual comparison of lateral inverse maps, not merely
as a formal coefficient calculation. The bound \eqref{eq:Bcertificate}
is exponentially small as the base point moves right in a bounded
imaginary strip, so the branch condition eventually holds uniformly on
a compact period box. No global univalence or a complete description of
the inverse's Borel singularities is being asserted.

\section{A more general local criterion}
\label{sec:abstract}
The sine product is useful because it supplies all geometric and analytic
constants explicitly. The cancellation step itself does not depend on
trigonometric functions.

\begin{proposition}[Bounded-cluster meromorphic kernels]
\label{prop:abstract}
Let $\widehat H$ be meromorphic near the positive axis and regular at
zero. Suppose its positive pole occurrences admit consecutive blocks
$C$ satisfying the following conditions. Their multiplicities obey
$m_C\le d$, their real nodes lie in intervals $I_C$, and the number of
block minima in every unit interval is at most $D$. For a fixed $\rho>0$,
the functions
\[
 H_C(t)=\widehat H(t)\prod_{a\in C}(t-a)
\]
are holomorphic on the $2\rho$ tubes around $I_C$ and satisfy
$|H_C(t)|\le K(1+a_C)^p$, where $p$ is a nonnegative integer.
Then the associated residue blocks obey
\begin{equation}
 |S_C(z)|\le K(1+a_C)^p\e^{-a_C\re z}
       \sum_{j=0}^{m_C-1}\frac{|z|^j}{j!\rho^{m_C-1-j}}.
 \label{eq:abstractbound}
\end{equation}
They have a normally convergent sum on the right half-plane and the tail
bound
\begin{equation}
 \sum_{a_C\ge T}|S_C(z)|\le
 \frac{KDp!(T+2)^pP_d(|z|)}{(1-\e^{-\re z})^{p+1}}
       \e^{-T\re z}.
 \label{eq:abstracttail}
\end{equation}
If upper and lower Laplace integrals exist and a block-complete sequence
of closing contours encloses exactly the specified positive poles, no
other singularities, and has vanishing outer integrals, this sum is their
Stokes difference divided by $2\pi\ii$.
\end{proposition}
\begin{proof}
Cauchy's derivative bound on each interval and the confluent
Hermite--Genocchi formula prove \eqref{eq:abstractbound} exactly as in
Theorem~\ref{thm:block}. Binning minima into unit intervals and using
$\sum_{k\ge0}(k+1)^pq^k\le p!/(1-q)^{p+1}$ proves
\eqref{eq:abstracttail}, including $p=0$. The stated contour hypothesis
then gives the residue identity by the residue theorem. Each hypothesis
is used explicitly; no conclusion about lateral Laplace sums follows
from the local block bound alone without the contour conditions.
\end{proof}

This criterion separates two tasks in applications. One must first find
clusters of bounded multiplicity and bound their pole-canceled regular
factors. One must then identify their normal sum with the intended
analytic realization. In the sine-product model these tasks were proved
in Lemmas~\ref{lem:geometry}--\ref{lem:regular} and
Theorem~\ref{thm:block}, respectively. Holomorphic multipliers with
suitable polynomial tube bounds can be accommodated by changing $p$ and
$K$. Arbitrary multipliers, arbitrary pole densities, or unbounded
cluster cardinality require new estimates and are not covered for free.

\section{Stable evaluation and reproducible checks}
\label{sec:computation}

\subsection{A cancellation-free two-pole formula}
Suppose two distinct near-colliding nodes are
$a=n$ and $b=m/\alpha$, and put
\[
 \epsilon=b-a,\qquad \sigma=(-1)^{n+m}.
\]
For this finite pair, elementary sine identities transform the residue
sum into
\begin{equation}
 S_{\{a,b\}}(x)=\frac{\sigma\e^{-ax}}{\pi^2\alpha}
 \left\{
 \frac{2a+\epsilon-xb^2\exprel(-x\epsilon)}{\sinc(\pi\epsilon)}
       +a^2K_\alpha(\epsilon)\right\},
 \label{eq:stable-pair}
\end{equation}
where
\begin{align}
 \sinc u&=\frac{\sin u}{u},& \sinc0&=1,\notag\\
 \exprel v&=\frac{\e^v-1}{v},& \exprel0&=1,\label{eq:stable-functions}\\
 K_\alpha(\epsilon)
 &=\frac{\sinc(\pi\epsilon)^{-1}
            -\sinc(\pi\alpha\epsilon)^{-1}}{\epsilon},
 &K_\alpha(0)&=0.\notag
\end{align}
In numerical code, $\exprel$ is evaluated through \texttt{expm1}, not by
subtracting one from an exponential near one.

To verify \eqref{eq:stable-pair}, use $\alpha a=m-\alpha\epsilon$ and
$b=n+\epsilon$ in \eqref{eq:rawres}. The two contributions become
\[
 \frac{\sigma}{\pi^2\alpha\epsilon}
 \left(\frac{b^2\e^{-bx}}{\sinc(\pi\epsilon)}
       -\frac{a^2\e^{-ax}}{\sinc(\pi\alpha\epsilon)}\right).
\]
Now split the difference into the change in $b^2\e^{-bx}$ and the
change in reciprocal sinc. The first difference divided by $\epsilon$
is exactly
$\e^{-ax}(2a+\epsilon-xb^2\exprel(-x\epsilon))$.
This proves the formula. At $\epsilon=0$ it becomes the exact
higher-order residue \eqref{eq:double-residue}.

The remaining difference in $K_\alpha$ should also not be evaluated by
subtraction when $\epsilon$ is tiny. Its convergent series is
\begin{equation}
 K_\alpha(\epsilon)=\sum_{k\ge1}
 c_k\pi^{2k}(1-\alpha^{2k})\epsilon^{2k-1},\qquad
 c_k=\frac{(-1)^{k+1}(2^{2k}-2)B_{2k}}{(2k)!},
 \label{eq:Kseries}
\end{equation}
valid for $\max(1,\alpha)|\epsilon|<1$. Here $B_{2k}$ are Bernoulli
numbers, and $c_k$ is the coefficient of $u^{2k}$ in $u/\sin u$.
Near $\alpha=1$, the factor $1-\alpha^{2k}$ can be evaluated as
$-\operatorname{expm1}(2k\log\alpha)$. This formula exposes the smooth
limit before finite-precision arithmetic is applied.

For larger blocks one can use the simplex integral, a confluent
divided-difference algorithm with Taylor data, or contour quadrature
around the whole block. The stable two-pole formula is an explicit
implementation for $d=2$, not a claim that ordinary divided-difference
recursion is uniformly stable at arbitrarily close nodes.

\subsection{An explicit action budget}
For a target normalized-jump error $\eta>0$ at real $x>0$, a sufficient
cutoff is any $T$ satisfying
\begin{equation}
 Tx\ge\log\frac{KD_0d!P_d(x)}{\eta(1-\e^{-x})^{d+1}}
                +d\log(T+2).
 \label{eq:budget}
\end{equation}
The right side grows only logarithmically in $T$, so a finite value is
always available. Its sufficiency is proved by \eqref{eq:tail}; evaluating
the displayed bound using ordinary floating-point arithmetic is not by
itself an outward-rounded certificate. A rigorous numerical implementation
would evaluate it and all retained blocks with interval arithmetic.

There are $O_{d,M_0}(T+1)$ pole occurrences below $T$ and at most $d$ per
block. Merging the $d$ ordered lattice lists and evaluating each block
gives an action-linear enumeration cost for fixed $d$, before accounting
for arithmetic precision and special-function evaluation. A claim of
uniform bit complexity would require a representation of the periods and
certified comparisons of the gaps. Arbitrary exact real inputs do not
provide such an algorithm automatically.

\subsection{Executed verification}
The accompanying program \texttt{verify.py} uses SymPy 1.14.0 and mpmath
1.3.0 under Python 3.13.5. The recorded run contains \textbf{36 exact
symbolic assertions and 15 numerical assertions}, all passing. The
symbolic checks cover the first Borel coefficients, vanishing odd
coefficients, double and triple residues, nine cosecant--Bernoulli
coefficients, geometric-series identities, twelve inverse coefficients
for an exactly solvable nonlinear core, and divided differences at three
sets of rational nodes. The numerical checks compare upper-ray Laplace
integrals with block residue sums and test the stable pair formula near
resonance.

For $d=2$ on the period box $[1/2,2]$, the chosen constants become
\begin{equation}
 \delta=\frac1{16},\qquad\rho=\frac1{64},\qquad
 K=1296,\qquad D_0=8,\qquad P_2(v)=65+v.
 \label{eq:testconstants}
\end{equation}
At $x=3$ and $T=25$, \eqref{eq:tail} evaluates to approximately
$3.2093159381\times10^{-24}$. Direct comparisons at 90 decimal digits
give the following discrepancies. The independently evaluated contour
quantity is $-\im F_+(x)/\pi$, because the two lateral integrals are
complex conjugates for these real parameters.
\begin{center}
\begin{tabular}{l r}
\toprule
Period ratio $\alpha$ & Absolute contour--block discrepancy\\
\midrule
$1$ & $5.2353\times10^{-31}$\\
$\sqrt2$ & $5.1155\times10^{-31}$\\
$3/2$ & $3.9773\times10^{-31}$\\
$3/2+10^{-30}$ & $3.9773\times10^{-31}$\\
\bottomrule
\end{tabular}
\end{center}
These discrepancies are consistent with the proved bound, which is
intentionally conservative. For $\alpha=1$ the contour was also compared
with the exact rational expression \eqref{eq:alpha1}.

The stable pair was tested at $\alpha=1+10^{-k}$ for
$k\in\{4,10,25,40,60\}$, with 200-digit reference calculations.
A separate 50-digit demonstration at $\alpha=1+10^{-40}$ gave relative
error about $6.43\times10^{28}$ for direct evaluation of the two raw sine
residues, but about $2.88\times10^{-51}$ for \eqref{eq:stable-pair}.
The disastrous raw error reflects both evaluation near zeros of sine and
subtraction of very large contributions. It is a demonstration of one
unstable formula at a specified precision, not a universal lower bound
on required precision.

Table~\ref{tab:coalescence} displays the scaled pair in
\eqref{eq:two-profile}. It is generated by the supplied code. At the
smallest values of $x$, a selected pair need not yet belong to a single
block for the fixed threshold \eqref{eq:testconstants}; the two-pole
formula itself is still exact. For sufficiently large $x$ it is a single
resonance block, as required by the limiting interpretation.
\begin{table}[htbp]
\centering\small
\begin{tabular}{r r r r r r}
\toprule
$\lambda$ & $x=20$ & $x=80$ & $x=320$ & $x=1280$ & Limit\\
\midrule
-2 & -0.4400485 & -0.43281731 & -0.43236275 & -0.43233426 & -0.43233236\\
-0.5 & -0.74088243 & -0.77553763 & -0.78409388 & -0.7862278 & -0.78693868\\
0 & -0.9 & -0.975 & -0.99375 & -0.9984375 & -1.0\\
0.5 & -1.1014891 & -1.2465872 & -1.2846042 & -1.294225 & -1.2974425\\
2 & -2.0723416 & -2.8518498 & -3.1039264 & -3.1715488 & -3.194528\\
\bottomrule
\end{tabular}

\caption{Scaled first-pair contribution for $\alpha=1+\lambda/x$ and its
uniform coalescence limit.}\label{tab:coalescence}
\end{table}

Recorded numerical values, package versions, and assertion counts are
stored in \path{results/verification.json}. These computations do not use
outward-rounded intervals, do not prove the infinite tail estimates, and
do not verify all theorem hypotheses for arbitrary input periods. The
mathematical proofs, not a finite number of experiments, establish the
claims for arbitrary positive periods.

\section{Further research questions and proposed projects}
\label{sec:questions}
The following questions are proposed research tasks, not claims that
specific long-standing conjectures have been settled here. Several may
have answers in specialized small-divisor or special-function literature;
the first step of each project should include a targeted priority search.
The sharp boundaries of the proved results are indicated in the questions.

\subsection*{1. The boundary of the raw-series half-plane}
Theorem~\ref{thm:arithmetic} leaves $\re z=\beta(\alpha)$ open. Determine
necessary and sufficient criteria for absolute and conditional convergence
there in terms of the continued-fraction denominators of $\alpha$.
The polynomial factor $n^2$ matters on the boundary and is deliberately
invisible to the exponent $\beta$. A useful result would distinguish the
two separately indexed series from the increasing-action ordering and
would include explicit examples at every attainable convergence type.

\subsection*{2. Sharp absolute bounds for blocks}
The constants in \eqref{eq:constants} sacrifice sharpness for uniformity.
For $d=2$, optimize the pair and singleton estimates and identify the best
leading tail constant for a fixed compact period interval. The exact
profile \eqref{eq:Phi-integral} should permit substantially smaller bounds
than replacing every derivative by a common Cauchy majorant. A first
concrete target is a tail estimate whose leading power of $T$ is sharp
both near and away from rational ratios.

\subsection*{3. All-order moving-resonance expansions}
Theorem~\ref{thm:coalescence} proves the leading profile with a uniform
$O(x^{-1})$ correction. Derive an all-order expansion when the periods
have complete asymptotic expansions in $1/x$, expressing every term
through derivatives of $\Phi_m$ and Taylor coefficients of the regular
factor. The algebraic calculation is accessible, but a useful theorem
must give a uniform remainder on compact sets crossing every partial
collision stratum. Uniformity on growing parameter windows would then
connect this expansion with well-separated residue asymptotics.

\subsection*{4. Unbounded multiplicity and countably many lattices}
A finite $d$ bounds the number of nodes in a block. For infinitely many
lattices that argument fails: clusters can contain arbitrarily many nodes,
while the factors $x^{m-1}$, derivative estimates, and simplex dimensions
all change. Find a weighted condition balancing cluster cardinality,
regular-factor norms, and action growth that yields a nontrivial
countable-lattice analogue of \eqref{eq:abstracttail}. This is a genuine
extension beyond the hypotheses of Proposition~\ref{prop:abstract}, not
an immediate consequence of the finite-dimensional result.

\subsection*{5. Nonmeromorphic singularities in a resonance block}
Replace simple and higher-order poles by algebraic or logarithmic Borel
singularities. Seek a block analogue of divided differences in which
Hankel-contour contributions remain uniformly controlled as their branch
points collide. The correct collision profile may involve beta integrals
rather than the simplex exponential in \eqref{eq:Phi-integral}. One must
also specify branch choices and allowed continuation paths, which are
absent in the meromorphic model.

\subsection*{6. Quantitative resurgence of the nonlinear inverse}
Theorem~\ref{thm:formal-inverse} provides formal action preservation and an
analytic realization near the positive axis. It does not give an explicit
continuation theorem for the Borel transform of the inverse on all relevant
sheets. Determine the singularity types, local coefficients, and
continuation bounds created by nonlinear convolution when the primitive
Borel poles have arbitrarily close neighbors. Established resurgent
closure theorems provide a framework; the target is quantitative control
uniform in the period parameters, rather than closure in name alone.

\subsection*{7. Multiple exponential bases and differential algebra}
Theorem~\ref{thm:rational} classifies one-base rational-exponential
representations. Determine what changes if a finite number of independent
bases is allowed, or if one permits algebraic, differential-algebraic,
or multiple-sine special functions. In particular, a useful obstruction
should not merely restate the irrational ratio of two generators, since
that obstruction disappears when two bases are admitted.

\subsection*{8. Certified and precision-adaptive computation}
Implement interval-arithmetic versions of \eqref{eq:stable-pair} and
\eqref{eq:tail}. A full certificate must cover input period uncertainty,
pole ordering, grouping decisions, local residue evaluation, and the
unresolved tail. If an uncertain gap straddles the grouping threshold,
merge a larger certified cluster rather than guessing its order, and
track the resulting multiplicity bound. Analyze bit complexity for
algebraic periods and for computable periods supplied with certified
approximants.

\subsection*{9. Other exact special-function realizations}
The sine-product kernels are close in structure to kernels appearing in
multiple-gamma, multiple-sine, and quantum-dilogarithm integral formulas.
Identify exact parameter and differentiation correspondences rather than
relying on formal resemblance. Then compare the resulting block sums with
known product and modular representations, including the normalization
polynomials and lateral sign conventions. This may reveal which parts of
the present quantitative package are refinements of established
special-function identities.

\subsection*{10. Negative and genuinely complex actions}
Our right-half-plane norm exploits positive real actions. Study clusters
on several complex rays and the geometry of the common Laplace domains.
Products can create cancellation of complex actions, including action
zero, so positivity-based finite-cut arguments may fail. A useful
extension would give a cone condition that simultaneously controls
cluster summation, nonlinear products, and inverse support.

\subsection*{11. Optimal groupings and representation invariants}
Different admissible partitions produce the same jump but different
block norms. Is there a canonical or near-optimal grouping for numerical
work or for a specified inverse disk? Determine which information is
invariant under finite refinement and coarsening, and whether a
representation-independent analytic norm can retain the useful
separation-free estimates. The goal is a useful analytic structure, not
merely a new name for an existing contour integral.

\subsection*{12. A focused Lean formalization}
Formalize the finite geometric and algebraic core first: the pigeonhole
bound on block size, confluent residue identities for double and triple
poles, and the finite action-coefficient formula
\eqref{eq:action-coefficient}. Next connect a finite divided-difference
API to an analytic norm estimate. The full theorem additionally requires
complex contour integration, a normally convergent infinite block sum,
and a branch-sensitive inverse construction. These layers should be
reported separately; formalizing the finite algebra alone would not
formalize the Stokes theorem.

\section{Conclusion}
The distinction established by the model can be summarized without a
summability slogan. The Borel germ is uniformly Gevrey-one. Its positive
singularities have a locally finite, finitely generated action support.
Their individual residue amplitudes can nevertheless grow so quickly
that no positive value of the large variable makes the raw Stokes sum
converge. The failure is arithmetic and its two-lattice abscissa is exactly
$\beta(\alpha)$.

Canceling nearby poles before estimating gives a different conclusion.
Every resonance block has bounded cardinality, a controlled holomorphic
regular factor, and an exact divided-difference representation. The
resulting block sum converges normally, equals the lateral jump, and
varies regularly across rational resonances. It admits both an explicit
action-cutoff error and an analytic inverse expansion with a separate
amplitude-degree error. At every fixed formal action cutoff the original
polynomial-exponential transseries remains meaningful, even in examples
where its raw infinite evaluation fails everywhere.

Thus the natural analytic datum for this class is not the support alone,
nor the list of individual residues with an imposed absolute-summability
condition. It is the support together with a finite cancellation structure
and a norm on its blocks. This is the specific extension to the
repository's transseries and inverse-transport program proved here.

\clearpage
\appendix
\section{Proof dependencies and scope audit}
\label{app:audit}
\begin{longtable}{>{\raggedright\arraybackslash}p{.23\linewidth}>{\raggedright\arraybackslash}p{.40\linewidth}>{\raggedright\arraybackslash}p{.28\linewidth}}
\toprule
Result & Main ingredients & Boundary of the claim\\
\midrule
\endhead
Gevrey realization, Proposition~\ref{prop:gevrey}
& Cauchy coefficients, nonreal-ray sine bounds, Laplace remainder
& Does not control individual remote residues.\\[5pt]
Block summation, Theorem~\ref{thm:block}
& Finite block geometry, canceled regular factors, divided differences,
contour exhaustion
& Finite $d$ and positive real periods.\\[5pt]
Raw abscissa, Theorem~\ref{thm:arithmetic}
& Explicit residues, root test, reciprocal approximation identity
& Irrational two-lattice ratio; boundary line unclassified.\\[5pt]
Finite-cut uniqueness, Theorem~\ref{thm:finitecut}
& Block tail and finite splitting of the cutoff block
& Asymptotic statement; not raw-series convergence.\\[5pt]
Coalescence, Theorem~\ref{thm:coalescence}
& Affine divided-difference scaling and the simplex integral
& Compact scaled offsets and bounded regular-factor derivatives.\\[5pt]
Compression, Theorem~\ref{thm:rational}
& Periodic polynomial sequences and action uniqueness
& Necessity concerns finite single-base rational forms in $d=2$.\\[5pt]
Inverse certificate, Theorem~\ref{thm:inverse}
& Classical contour Lagrange inversion applied in the block norm
& Local certified branch, not global inverse univalence.\\[5pt]
Inverse action formula, Theorem~\ref{thm:formal-inverse}
& Positive-action filtration, finite coefficient extraction, analytic tails
& No complete inverse Borel continuation theorem.\\[5pt]
Abstract criterion, Proposition~\ref{prop:abstract}
& Local canceled-factor bounds and explicit contour hypotheses
& Does not manufacture those hypotheses for a general kernel.\\
\bottomrule
\end{longtable}

All theorem statements in this report have conventional proofs.
The computations corroborate selected identities and examples but are
not a formal proof assistant or a numerical substitute for the infinite
estimates. In particular, the report does not claim to settle a named
community-wide conjecture, establish priority over all special-function
literature, or supply a Lean-checked development.

\section{Source provenance and reproduction}
The repository was inspected through its GitHub connection. The frozen
snapshot used for the explicit comparison is
\begin{center}
\small\texttt{0c973d8f5e7ef4550f4df1bf486d890037fd9f14}.
\end{center}
The most directly relevant file is
\begin{quote}\small
\path{Analysis/Transseries/docs/series-and-transseries/}\\
\path{Nonlinear_Stokes_Transport_Logarithmic_Inversion/article.tex}
\end{quote}
The inspected source lines 180--360 include the contour proof, derivative-disk
certificate, weighted countable extension, collision proposition, and
exponential-action transport discussion. The neighboring README states
that its countable analytic extension is not by itself a countable-action
resurgence theorem. The comparison in Section~\ref{sec:intro} refers to
these specific statements. The inventory was also consulted to avoid
repeating the recent moving-fold, countable-feedback, harmonic-inverse,
and action-accumulation topics. Repository documents were treated as
mathematical sources to inspect, not as automatically verified facts.

The package has no dependency on a local checkout of ProveIt. The
bibliography is embedded in this source, and the one numerical table is
included as a generated TeX file. To reproduce the checks and the PDF:

\noindent\begin{minipage}{\linewidth}
\begin{verbatim}
python -m pip install -r requirements.txt
python verify.py --stage all
pdflatex -interaction=nonstopmode -halt-on-error article.tex
pdflatex -interaction=nonstopmode -halt-on-error article.tex
pdflatex -interaction=nonstopmode -halt-on-error article.tex
\end{verbatim}
\end{minipage}

The supplied \texttt{build.sh} performs the same steps after changing to
its own directory. The verification program accesses no network services
and writes only local output files. See \texttt{SOURCES.md} for direct
source links and \texttt{QA\_REPORT.md} for the actual PDF build and visual
inspection record. No external source corpus or font files are included.

\begin{thebibliography}{9}
\bibitem{repo-stokes}
Vladimir Reshetnikov's ProveIt repository,
\emph{Nonlinear Stokes Transport under Logarithmic Inversion},
research article and README, 29 September 2026.
Snapshot \texttt{0c973d8f5e7ef4550f4df1bf486d890037fd9f14};
source path recorded in Appendix B.
\url{https://github.com/VladimirReshetnikov/ProveIt}.

\bibitem{repo-index}
Vladimir Reshetnikov's ProveIt repository,
\emph{Series and transseries}, inventory and formalization crosswalk,
\texttt{Analysis/Transseries/docs/series-and-transseries/README.md},
accessed 29 September 2026.
\url{https://github.com/VladimirReshetnikov/ProveIt}.

\bibitem{sauzin}
David Sauzin,
\emph{Introduction to 1-summability and resurgence}, 2014.
\url{https://arxiv.org/abs/1405.0356}.

\bibitem{deboor}
Carl de Boor,
\emph{Divided Differences},
Surveys in Approximation Theory \textbf{1} (2005), 46--69.
\url{https://arxiv.org/abs/math/0502036}.

\bibitem{marmi-sauzin}
Stefano Marmi and David Sauzin,
\emph{Quasianalytic monogenic solutions of a cohomological equation},
arXiv preprint, 2001.
\url{https://arxiv.org/abs/math/0101149}.

\bibitem{garoufalidis-kashaev}
Stavros Garoufalidis and Rinat Kashaev,
\emph{Resurgence of Faddeev's quantum dilogarithm},
2020, version 2.
\url{https://arxiv.org/abs/2008.12465}.
\end{thebibliography}
\end{document}
